\documentclass[12pt]{article}

% setting margins for the document
\usepackage[letterpaper, margin=1.1in]{geometry}

\usepackage[T1]{fontenc}
% packages used for the various math commands
\usepackage{amsmath, amsthm, amssymb,enumitem}

% setting up the font
\usepackage{libertinus-type1}
\usepackage{libertinust1math}

% adjusting in-word spacing
\usepackage{microtype}

\usepackage{parskip}
\usepackage{etoolbox}

% setting custom margins for enumerate and itemizaed lists
\setlist[enumerate]{parsep=0pt plus 4pt,topsep=3pt plus 4pt}
\setlist[itemize]{parsep=0pt plus 4pt,topsep=3pt plus 4pt,itemsep=.5ex}

\leftmargini=5.5ex
\leftmarginii=3.5ex


%For separated lists with consecutive numbering
\newcounter{separated}
\newcommand{\excise}[1]{}
\newcommand{\comment}[1]{{$\star$\sf\textbf{#1}$\star$}}
\newcommand{\exnumber}[1]{\noindent\textbf{#1}}

%new math symbols taking no arguments
\newcommand\0{\mathbf{0}}
\newcommand\CC{\mathbb{C}}
\newcommand\FF{\mathbb{F}}
\newcommand\NN{\mathbb{N}}
\newcommand\QQ{\mathbb{Q}}
\newcommand\RR{\mathbb{R}}
\newcommand\ZZ{\mathbb{Z}}
\newcommand\bb{\mathbf{b}}
\newcommand\kk{\Bbbk}
\newcommand\mm{\mathfrak{m}}
\newcommand\xx{\mathbf{x}}
\newcommand\yy{\mathbf{y}}
\newcommand\GL{\mathit{GL}}
\newcommand\pset{\mathcal{P}}
\newcommand\into{\hookrightarrow}
\newcommand\nsub{\trianglelefteq}
\newcommand\onto{\twoheadrightarrow}
\newcommand\minus{\smallsetminus}
\newcommand\goesto{\rightsquigarrow}

% theorem environment
\newtheorem*{theorem}{Theorem}

%redefined math symbols taking no arguments
\newcommand\<{\langle}
\renewcommand\>{\rangle}
\renewcommand\iff{\Leftrightarrow}
\renewcommand\phi{\varphi}
\renewcommand\implies{\Rightarrow}

%new math symbols taking arguments
\newcommand\ol[1]{{\overline{#1}}}

%redefined math symbols taking arguments
\renewcommand\mod[1]{\ (\mathrm{mod}\ #1)}

%roman font math operators
\DeclareMathOperator\im{im}

%for easy 2 x 2 matrices
\newcommand\twobytwo[1]{\left[\begin{array}{@{}cc@{}}#1\end{array}\right]}

%for easy column vectors of size 2
\newcommand\tworow[1]{\left[\begin{array}{@{}c@{}}#1\end{array}\right]}

% redefining \maketitle so that it has proper spacing
\makeatletter
\renewcommand{\maketitle}{%
  \begingroup
  \setlength{\parskip}{0pt}   % kills the parskip effect inside the title block
  \begin{center}
    {\LARGE\bfseries \@title\par}
    \vspace{1.5ex}
    {\large \@author\par}
    \vspace{0.75ex}
    {\normalsize \@date\par}
  \end{center}
  \endgroup
  \vspace{2ex}
}
\makeatother

% makes sure that proofs within an itemized or enumerate environment 
% maintain proper spacing between problem statements. 
\makeatletter
\BeforeBeginEnvironment{proof}{%
  \par
  \ifnum\@listdepth>0 \addvspace{\medskipamount}\fi
}
\AfterEndEnvironment{proof}{%
  \ifnum\@listdepth>0 \addvspace{\medskipamount}\fi
}
\makeatother

%%%%%%%%%%%%%%%%%%%%%%%%%%%%%%%%%%%%%%%%%%%%%%%%%%%%%%%%%%%%%%%%%%%%%%
\begin{document}%%%%%%%%%%%%%%%%%%%%%%%%%%%%%%%%%%%%%%%%%%%%%%%%%%%%%%
%%%%%%%%%%%%%%%%%%%%%%%%%%%%%%%%%%%%%%%%%%%%%%%%%%%%%%%%%%%%%%%%%%%%%%
	\title{Finite Sets Solutions}
	\author{Solutions by: Raindrops\_drop}
	\date{October 11, 2026}
	\maketitle
	
	\section*{Exercises}

    \exnumber{1.} 
    \begin{enumerate}[label=(\alph*)]
        \item Make a list of all the injective maps 
        $$
        f: \{1, 2, 3\} \to \{1, 2, 3, 4\}.
        $$
        Show that none is bijective. (This constitutes a \textit{direct}
        proof that a set $A$ of cardinality three does not have cardinality
        four).
        \begin{proof}
            Define $f_1$ such that $f_1(i) = i$, note this is not surjective since
            no $i$ will map onto 4.

            Define $f_2$ such that $f_2(i) = i + 1$, note this is not surjective 
            since no $i$ will map onto 1.

            Define $f_3$ such that $f_3(1) = 1$, and $f_3(i) = i + 1$ for 
            the other $i$. Note this is not surjective since no $i$ will map
            onto 2.

            Define $f_4$ such that $f_4(1) = 1$, $f_4(2) = 2$, and 
            $f_4(3) = 4$. Note this is not surjective since no $i$ will map onto 3.

            For each $f_i$ we can construct 5 other injective functions by cycling
            what element each one gets mapped onto. Note, however, that this
            cycling process does not change what element in the co-domain 
            remains not covered by $f_i$. As such, we conclude that 
            there exists no bijective correspondence between these two sets.
        \end{proof}
        \item How many injective maps 
        $$
        f: \{1, \ldots, 8\} \to \{1, \ldots, 10\}
        $$
        are there? (You can see why one would not wish to try to prove 
        directly that there is no bijective correspondence between these
        sets).
        \begin{proof}
            Note that under an injective map $f$, the element 
            $1$ has 10 possible choices on what to get mapped onto, the 
            element 2 then has 9 possible choices on what to get mapped onto.
            Following this logic, we note that there are 
            $10 \cdot 9 \cdot 8 \cdot 7 \cdot 6 \cdot 5 \cdot 4 \cdot 3$
            possible injective maps, or simply $10!/2$ which is equivalent 
            to 1,814,400.
        \end{proof}
    \end{enumerate}

    \exnumber{2.} Show that if $B$ is not finite and $B \subseteq A$,
    then $A$ is not finite.

    \begin{proof}
        Let $B \subseteq A$ such that $B$ is not finite. Assume for contradiction
        $A$ is finite. Then, there exists an injective function 
        $f: A \into S_n$ where $S_n$ is a section of the positive integers.
        Since $B \subseteq A$, note the restriction of $f$, $g: B \into S_n$
        is also injective. This implies $B$ is finite, which is a contradiction.
        Therefore, it must be $A$ is not finite. 
    \end{proof}

    \exnumber{3.} Let $X$ be the two-element set $\{0, 1\}$. Find a bijective
    correspondence between $X^\omega$ and a proper subset of itself. 

    \begin{proof}
        \comment{Section not finished}
    \end{proof}

    \exnumber{4.} Let $A$ be a non-empty finite simply ordered set.
    \begin{enumerate}[label=(\alph*)]
        \item Show that $A$ has a largest element. [Hint: Proceed by induction
        on the cardinality of $A$].
        \begin{proof}
            If $A$ has cardinality $1$, then $A$ has only one element,
            and thus said element is also the largest element. Assume now
            for induction that a non-empty finite simply ordered set of cardinality $n$ has a largest
            element. Now consider $A$ with cardinality $n + 1$. 
            Let $a \in A$ and define $A_n = A - \{a\}$, noting $A_n$ has
            cardinality $n$. By our induction assumption, $A_n$ has a largest
            element, call it $a_n$. If $a_n > a$ then $a_n$ is the largest
            element of $A$, otherwise if $a_n < a$, then since $a_n$
            is the largest element of $A_n$, we have $a_n \geq p$ for all
            $p \in A_n$, implying thus $a > p$ for all $p \in A_n$.
            As such, $a$ is the largest element of $A$. In either case 
            one can conclude $A$ has a largest element.
        \end{proof}
        \item Show that $A$ has the order type of a section of the positive
        integers.
        \begin{proof}
            Since $A$ is finite, there exists a bijection from $A$ 
            onto a section of the positive integers $S_n$.
            If the cardinality of $A$ is one, then the single existing
            bijection is trivially order preserving.
            Assume for induction that for a non-empty finite simply order set $A$
            of cardinality $n$, it has the same order type as $S_n$, i.e.,
            there exists an order preserving bijection between the two sets.
            Consider now the case $A$ has cardinality $n + 1$. 
            Let $a \in A$ be the largest element of $A$ (which exists by part (a)), 
            and consider the subset $A_n \subset A$
            defined as $A_n = A - \{a\}$. Note by our induction assumption
            there exists an order preserving bijection $g: A_n \to S_n$.
            Define now $f: A \to S_{n+1}$ as $f(p) = g(p)$ for $p \in A_n$
            and $f(a) = n + 1$. Note $f$ is order preserving on all $p \in A_n$
            and since $a$ and $n + 1$ are the largest element of $A$ and
            $S_n$ respectively, $f$ is also order preserving with respect to
            them. Thus, $f$ is an order preserving bijection between
            $A$ and $S_{n+1}$.
        \end{proof} 
    \end{enumerate}

    \exnumber{5.} If $A \times B$ is finite, does it follow that $A$ and $B$
    are finite?

    \begin{proof}
        It does not follow from $A \times B$ that $A$ and $B$ are finite.
        Consider the sets $A = \RR^\omega$ and $B = \emptyset$, then
        $A \times B = \emptyset$. Thus, we noted an example of a 
        finite set created by the Cartesian product of an infinite set
        and the empty set. 
    \end{proof}

    \exnumber{6.} 
    \begin{enumerate}[label=(\alph*)]
        \item Let $A = \{1, \ldots, n\}$. Show there is a bijection of 
        $\pset(A)$ with the Cartesian product $X^n$, where $X$ is the 
        two-element set $X = \{0, 1\}$.
        \begin{proof}
            Define the function $f: \pset(A) \to X^n$ by 
            the map $S \subseteq A$ is mapped to $(x_1, x_2, \ldots, x_n)$
            where $x_i = 1$ if $i \in S$ and $x_i = 0$ otherwise,
            noting $1 \leq i \leq n$. To show that this function is 
            injective, assume $S_1, S_2 \subseteq A$ such that
            $f$ maps both subsets of $A$ onto the same element of 
            $X^n$. Then, if $i \in S_1$, it implies $x_i = 1$,
            also making it that $i \in S_2$. Thus, $S_1 \subseteq S_2$.
            By an analogous argument we have $S_2 \subseteq S_1$, and 
            thus $S_1 = S_2$. Hence, $f$ is injective.
            
            To show $f$ is surjective, let $(x_1, \ldots, x_n) \in X^n$
            be an arbitrary element of $X^n$. Let $I$ be the set contain all the indexes
            $i$ such that $x_i = 1$. Note then $I \subseteq A$, and thus
            $I \in \pset(A)$. In particular, $f(I) = (x_1, \ldots, x_n)$
            by definition of $f$. Therefore, $f$ is also surjective, and we 
            thus conclude $f$ is a bijection. 
        \end{proof}
        \item Show that if $A$ is finite, then $\pset(A)$ is finite.
        \begin{proof}
            It was proved in an earlier section that if a set has $n$ 
            elements, its power set contains $2^n$ elements. 
            Thus, if $A$ is finite and has $n$ elements, then
            $\pset(A)$ has $2^n$ elements. In particular, this
            shows $\pset(A)$ is finite if $A$ is finite.
        \end{proof}
    \end{enumerate}

    \exnumber{7.} If $A$ and $B$ are finite, show that the set of 
    functions $f: A \to B$ is finite.
    \begin{proof}
        Let $A$ and $B$ be finite sets. Note that any function
        $f: A \to B$ must map every element of $A$ to some element
        in $B$. As such, if $A$ has cardinality $n$ and $B$ has cardinality
        $m$, each element of $A$ has $m$ possible elements it can be 
        mapped onto. Thus, there must be $m^n$ total 
        distinct functions that map $A$ to $B$. As such,
        there is a finite total amount of functions $f: A \to B$.
    \end{proof}
%%%%%%%%%%%%%%%%%%%%%%%%%%%%%%%%%%%%%%%%%%%%%%%%%%%%%%%%%%%%%%%%%%%%%%
\end{document}%%%%%%%%%%%%%%%%%%%%%%%%%%%%%%%%%%%%%%%%%%%%%%%%%%%%%%%%
%%%%%%%%%%%%%%%%%%%%%%%%%%%%%%%%%%%%%%%%%%%%%%%%%%%%%%%%%%%%%%%%%%%%%%
